\PuzzleChapterIntro{Two Clues, One Suspect}{Bayes \textbullet\ Induction \textbullet\ Invariants \textbullet\ Sums of Squares}{A whodunit built from two dependent-looking “clues” whose numerical weights are themselves forged through three gates: an extremal interval invariant (frames), an induction-counting lemma (crossings), and a sum-of-two-squares certification via modular arithmetic (prime dial). Expect to translate structure into probabilities, handle a hidden mixture over imitation, and then fuse the evidence cleanly with Bayes to extract a single negative-integer score.}

\Lore{
The Archivists call it \textit{the Walk of Three Gates}.

Frames first: one wrong shadow spoils the rest.
Crossings next: safe paths are counted, never guessed.
Last, the Prime Dial: it turns only when a prime agrees to be two squares.

A bowl of coins waits on the Crucible table, bisected by a single scratched line.
The Initiate is told to ignore the noise, and compute what remains after both coins fall.
}

\Problem
\small
Two suspects are considered for the theft: \textbf{Eld} and \textbf{Nyra}.
\[
P(\text{Eld})=\frac25,\qquad P(\text{Nyra})=\frac35.
\]
If Nyra is guilty, she may mirror Eld's style. Let \(H=\) “Nyra mirrors”, with
\[
P(H\mid\text{Nyra})=\frac12.
\]

\medskip
\noindent\textbf{Gate I (frames).}
A \emph{frame} is an axis-aligned rectangle with positive side lengths.
A set \(B_1,\dots,B_F\) is arranged so that
\[
B_i\cap B_j\neq\varnothing \iff i\not\equiv j\pm1\pmod F,
\]
and the thief always uses as many frames as possible (so \(F\) is maximal).
The \textbf{Frame Coin} is flipped by choosing an unordered pair of distinct frames uniformly;
let \(E=\) “the chosen pair is disjoint”.
Eld always uses the cyclic rule above. If Nyra is guilty: if \(H\), she uses the same rule; if \(\neg H\),
she uses a \emph{chain} rule in which the disjoint pairs are exactly
\(\{(B_1,B_2),\dots,(B_{F-1},B_F)\}\).

\medskip
\noindent\textbf{Gate II (crossings).}
Using the same \(F\), draw \(F\) lines in the plane, no two parallel and no three concurrent,
and choose a point \(O\) on no line.
A crossing \(X\) (intersection of two lines) is \emph{clear} if the open segment \((OX)\) meets at most \(F-3\) lines.
Let \(C\) be the minimum number of clear crossings that must exist (this minimum is known to be sharp).

\medskip
\noindent\textbf{Gate III (prime dial).}
Set \(P=4C+1\) (for the value of \(C\) above, \(P\) is prime), and define
\[
W\equiv (2C)!\pmod P,\qquad 1\le W\le \frac{P-1}{2},\qquad
Y=\frac{W^2+1}{P}\in\mathbb Z_{>0}.
\]
The Prime Gate certifies that there exist nonnegative integers \(g<h\) and \(m,n\) such that
\[
P=g^2+h^2,\qquad Y=m^2+n^2,\qquad W=gm+hn.
\]
The \textbf{Prime Coin} is flipped by choosing one of the \(P\) dial detents uniformly;
let \(J=\) “the detent is dangerous”, with
\[
P(J\mid\text{Eld})=\frac{h}{P},\quad
P(J\mid\text{Nyra},\neg H)=\frac{g}{P},\quad
P(J\mid\text{Nyra},H)=\frac{g+h}{P}.
\]

\medskip
\noindent\textbf{Dependence.}
Assume \(E\) and \(J\) are independent conditional on \(\text{Eld}\), and also independent conditional on
\((\text{Nyra},H)\) and \((\text{Nyra},\neg H)\).

\medskip
You observe \emph{both} \(E\) and \(J\).

\medskip
\VaultDirective{Let $P(\text{Eld}\mid E,J)=\frac{A}{B}$ in lowest terms. Compute $A-B$.}

\SolutionPage

\scriptsize

\textbf{Roadmap.} Write \(EJ:=E\cap J\). We compute \(F\to C\to(P,W,Y)\to(g,h)\), then the likelihoods
\(P(EJ\mid\text{Eld})\) and \(P(EJ\mid\text{Nyra})\) (Nyra is a \(\tfrac12/\tfrac12\) mixture over \(H\)),
and finish with Bayes.

\medskip
\textbf{Gate I: find \(F\).}
Let \(X_i=[a_i,b_i]=\mathrm{proj}_x(B_i)\) and \(Y_i=[c_i,d_i]=\mathrm{proj}_y(B_i)\).
If \(B_i\cap B_j\neq\varnothing\) then \(X_i\cap X_j\neq\varnothing\) and \(Y_i\cap Y_j\neq\varnothing\); hence
\[
B_i\cap B_{i+1}=\varnothing \ \Rightarrow\ (X_i\cap X_{i+1}=\varnothing)\ \text{or}\ (Y_i\cap Y_{i+1}=\varnothing).
\tag{$\ast$}
\]
Choose indices so that \(a_1=\max\{a_1,\dots,a_F\}\). Then for every \(i\not\equiv 2,F\pmod F\),
the rule forces \(B_i\cap B_1\neq\varnothing\), hence \(a_1\in X_i\); in particular
\(X_i\cap X_{i+1}\neq\varnothing\) for \(3\le i\le F-2\).
Thus only the four cyclic pairs
\(X_1\cap X_2,\,X_2\cap X_3,\,X_{F-1}\cap X_F,\,X_F\cap X_1\) can be empty, and not all four can be
(empty \(X_1\cap X_2\) and \(X_2\cap X_3\) would force \(X_2\cap X_{F-1}=\varnothing\), contradicting the
required intersection \(B_2\cap B_{F-1}\neq\varnothing\)).
Hence \(\#\{i: X_i\cap X_{i+1}=\varnothing\}\le 3\), and similarly \(\#\{i: Y_i\cap Y_{i+1}=\varnothing\}\le 3\).
By \((\ast)\), each of the \(F\) disjoint adjacent frame-pairs is witnessed by an empty \(X\)-adjacency or \(Y\)-adjacency, so
$
F\le 3+3=6.
$
A configuration with \(F=6\) exists, so the maximal value is \(F=6\).

\medskip
\textbf{Gate II: find \(C\).}
Set \(t:=F-3\). The Crossing-Map lemma gives at least \(\frac12(t+1)(t+2)\) clear crossings, and this bound is sharp.
Induction spine: pick a line \(\ell\) closest to \(O\); choose \(P\in\ell\) so \((OP)\) meets no line.
Along each ray of \(\ell\) from \(P\), the number of lines crossed by \((OX)\) changes by at most \(1\) between consecutive
intersection points, yielding at least \(t+1\) clear crossings on \(\ell\); removing \(\ell\) reduces \(t\) by \(1\).
With \(F=6\), \(t=3\), so
$
C=\frac12(4)(5)=10.
$

\medskip
\textbf{Gate III: compute \(P,W,Y,g,h\).}
Now \(P=4C+1=41\) (prime) and \(W\equiv (2C)!\equiv 20!\pmod{41}\) with \(1\le W\le 20\).
By Wilson,
\[
40!\equiv -1\pmod{41},\qquad 40!=(20!)\,(21\cdots 40)\equiv (20!)^2\pmod{41},
\]
so \(W^2\equiv -1\pmod{41}\) and hence \(Y=(W^2+1)/41\in\mathbb Z_{>0}\).
Compute
\[
20!=\prod_{i=1}^{10} i(21-i)\equiv
20\cdot38\cdot13\cdot27\cdot39\cdot8\cdot16\cdot22\cdot26\cdot28 \pmod{41}.
\]
Grouping,
\[
(20\cdot39)(38\cdot27)(13\cdot28)(8\cdot16)(22\cdot26)\equiv
1\cdot1\cdot(-5)\cdot 5\cdot(-2)=50\equiv 9\pmod{41},
\]
so \(W=9\) and \(Y=(81+1)/41=2\). Now \(W^2+1=PY\) gives \(PY-W^2=1\).
Applying the \((xy-z^2=1)\Rightarrow\) “two-squares” lemma (the Prime Gate’s certificate),
there exist \(g<h\) and \(m,n\ge 0\) with
\[
P=g^2+h^2,\qquad Y=m^2+n^2,\qquad W=gm+hn.
\]
Since \(Y=2\), we have \((m,n)=(1,1)\), so \(W=g+h\).
With \(g^2+h^2=41\) and \(g+h=9\), we get \(gh=20\), hence \((g,h)=(4,5)\).

\medskip
\textbf{Likelihoods: }
With \(F=6\),
$
P(E\mid\text{cyclic})=\frac{F}{\binom{F}{2}}=\frac{2}{F-1}=\frac25,\qquad
P(E\mid\text{chain})=\frac{F-1}{\binom{F}{2}}=\frac{2}{F}=\frac13.
$


Thus
$
P(E\mid\text{Eld})=\frac25,\quad P(E\mid\text{Nyra},H)=\frac25,\quad P(E\mid\text{Nyra},\neg H)=\frac13,
$
and
\[
P(J\mid\text{Eld})=\frac{h}{P}=\frac{5}{41},\quad
P(J\mid\text{Nyra},\neg H)=\frac{g}{P}=\frac{4}{41},\quad
P(J\mid\text{Nyra},H)=\frac{g+h}{P}=\frac{9}{41}.
\]
Independence gives
\[
P(EJ\mid\text{Eld})=\frac25\cdot\frac{5}{41}=\frac{2}{41},
\]
\[
P(EJ\mid\text{Nyra})=\frac12\!\left(\frac25\cdot\frac{9}{41}\right)+\frac12\!\left(\frac13\cdot\frac{4}{41}\right)
=\frac{9}{205}+\frac{2}{123}=\frac{37}{615}.
\]

\medskip
\textbf{Bayes: }
\[
P(\text{Eld}\mid E,J)=
\frac{\frac25\cdot\frac{2}{41}}{\frac25\cdot\frac{2}{41}+\frac35\cdot\frac{37}{615}}
=\frac{\frac{4}{205}}{\frac{4}{205}+\frac{111}{3075}}
=\frac{60}{60+111}
=\frac{20}{57}.
\]
So \(A=20\), \(B=57\), hence \(A-B=-37\).

\FinalAnswer{\(\displaystyle -37\)}

\NextPage
